% !TeX root = ../navier-stokes-manuscript.tex
\clearpage
\section{The Missing Estimate and Global Smoothness}\label{sec:TheMissingEstimateAndGlobalSmoothness}

The Tower has reached the place where a smooth solution can still disappear.
Every finite rectangle exists on every shorter interval, and every overlap
contains the same derivatives of the same fluid.  Neither fact carries those
norms to \(T_*\).  If the cutoff-dependent bounds rise without limit, the
intersection never reaches the terminal face and there is no datum from which
to restart.

This is why the order of the quantifiers in
\eqref{eq:BodyWholeTerminalRectangleEstimate} matters.  We fix one finite
rectangle first and then move the cutoff.  Its bound may know \(u_0\), \(\nu\),
\(T_*\), \(N\), and \(M\).  It cannot know how close \(T\) is to \(T_*\).
Once that estimate is paid, the compatible rectangles become one mixed
intersection, the intersection reaches one smooth terminal datum, and the same
Navier--Stokes history restarts from it.  We follow that complete argument
before returning to the critical height, the Zeno ladder, and the narrowing
vortex that forced us to build the Tower.

\subsection{The intersection reaches one terminal datum}\label{sub:NoFiniteSingularTime}

\begin{theorem}[Whole-space existence and smoothness]\label{thm:MainGlobalRegularity}
For every \(\nu>0\) and every \(u_0\in\Ssigma\), the maximal classical solution of \eqref{eq:NavierStokesSystem} satisfies
\begin{equation}\label{eq:BodyGlobalSmoothnessConclusion}
  T_*=\infty,
  \qquad
  u,p\in C^\infty\!\bigl(\R^3\times[0,\infty)\bigr).
\end{equation}
For every \(t\geq0\),
\begin{equation}\label{eq:GlobalEnergyConclusion}
  \frac12\norm{u(t)}_2^2
  +\nu\int_0^t\norm{\nabla u(s)}_2^2\,ds
  =\frac12\norm{u_0}_2^2.
\end{equation}
\end{theorem}

\begin{proof}
Suppose the maximal history ends at some finite \(T_*\).
For each finite pair \((N,M)\),
\Cref{thm:BodyWholeTerminalRectangleTheorem} carries its pressure-complete
rectangle all the way across \(I_*\).  Whenever two rectangles overlap,
\eqref{eq:BodyRectangleRestrictionCompatibility} says that they contain the
same derivatives of the same history.  Their compatible union gives
\[
  u\in\bigcap_{r,m<\infty}H^r(I_*;H^m(\R^3)).
\]

Now hold one spatial height \(m\) fixed.  The temporal row \(r=1\) puts
\(u\) in \(H^1(I_*;H^m)\), so the actual history has an \(H^m\) limit as
\(t\uparrow T_*\).  A higher spatial row cannot produce a different limit:
after embedding it into \(H^m\), both limits come from the same \(u(t)\).
All finite heights meet at one divergence-free endpoint
\begin{equation}\label{eq:BodyCommonSmoothEndpoint}
  u_*:=\lim_{t\uparrow T_*}u(t)
  \in\Hinfty_\sigma(\R^3)
  \subset\bigcap_{m<\infty}H^m(\R^3).
\end{equation}
This is the endpoint furnished by \Cref{lem:CompleteSpatialRowEndpoint}.

The terminal face must also remember how the history arrives there.  Given
finite \(N\) and \(m\), the row \(r=N+1\) places
\(V_N=\partial_t^Nu\) in \(H^1(I_*;H^m)\).  It has a trace \(V_N^*\), and
spatial compatibility puts that trace in \(\Hinfty\).  The normalized pressure
recurrence and the velocity recurrence survive the limit:
\begin{align*}
  Q_N^*
  &=R_iR_j\sum_{a=0}^{N}\binom{N}{a}V_{a,i}^*V_{N-a,j}^*,
  \\
  V_{N+1}^*
  &=\nu\Delta V_N^*
  -\sum_{a=0}^{N}\binom{N}{a}(V_a^*\cdot\nabla)V_{N-a}^*
  -\nabla Q_N^*.
\end{align*}
The traces are the velocity and Riesz-normalized pressure responses generated
by \(u_*\), as \Cref{prop:PressureCompleteEndpointJet} records.

Use \(u_*\) as initial data at time \(T_*\).  The smooth local construction in
\Cref{lem:SmoothLocalRestart} gives a right-hand history on
\([T_*,T_*+\delta]\).  Its recurrence begins from the same datum and the same
pressure normalization, so its entire jet agrees with the left-hand traces.
The join is smooth.  Strong uniqueness makes it the continuation of the
history generated by \(u_0\), contradicting the definition of \(T_*\).

Hence \(T_*=\infty\), and the normalized pair is smooth on every finite time
interval.  The kinetic-energy identity in
\Cref{prop:ClassicalKineticEnergyIdentity} now holds on every initial segment
of this global history and gives \eqref{eq:GlobalEnergyConclusion}.  In
particular,
\[
  \sup_{t\geq0}\int_{\R^3}|u(x,t)|^2\,dx
  \leq
  \int_{\R^3}|u_0(x)|^2\,dx
  <\infty.
\]
The rapid decay of \(u_0\) pays the last finiteness.  This is the smooth global
history with uniformly bounded kinetic energy required in
\eqref{eq:WholeSpaceRegularityTarget}.
\end{proof}

The theorem uses the whole intersection because that is the object the Tower
produced.  Its lowest rectangle gives a second, independent way to see why the
alleged last time has nowhere to form.

\subsection{The first rectangle seen from the endpoint}\label{sub:FromFiniteRectanglesToClassicalEstimates}

Choose the first rectangle, \((N,M)=(0,1)\).
Its adjacent rows are
\[
  u\in L^2(I_*;H^2),
  \qquad
  u_t\in L^2(I_*;L^2).
\]
Their trace lands in the space between them:
\[
  u\in C([0,T_*];H^1).
\]
Since \(H^1(\R^3)\hookrightarrow L^3(\R^3)\), this gives
\begin{equation}\label{eq:EndpointCriterionFromLowRectangle}
  \sup_{0\leq t\leq T_*}\norm{u(t)}_3<\infty.
\end{equation}
But the endpoint theorem of Escauriaza, Seregin, and \v{S}ver\'ak says that a finite maximal time forces the opposite behaviour \parencite{EscauriazaSereginSverak2003}:
\begin{equation}\label{eq:FiniteTimeLThreeEscape}
  T_*<\infty
  \quad\Longrightarrow\quad
  \sup_{0<t<T_*}\norm{u(t)}_3=\infty.
\end{equation}
The rectangle bounds the same velocity in the same norm on the same terminal interval.
The two conclusions cannot coexist, so this first coordinate already contradicts a finite \(T_*\).

The full intersection keeps every higher response available, but the physical
return does not need an inventory of them.  The first rectangle has already
reached the critical velocity norm.  The next one reaches the motion of the
critical height and the gradient that stretches the vortex.  Those are the two
coordinates we now follow.

\subsection{The critical-height roof}\label{sub:ClassicalContinuationCriteriaInsideFiniteRectangles}

Once the \(H^1\) trace reaches the terminal face, the critical height cannot keep
climbing underneath it.  Indeed,
\(H=\tfrac12\norm{\Lambda^{1/2}u}_2^2\), so
\begin{equation}\label{eq:BodyCriticalHeightRoof}
  H_{\max}
  :=\sup_{0\leq t\leq T_*}H(t)
  \leq\frac12\sup_{0\leq t\leq T_*}\norm{u(t)}_{H^1}^2
  <\infty.
\end{equation}
Start at any positive height \(H_0\), double it successively by
\(L_k:=2^kH_0\), and count the levels the same history actually reaches:
\[
  N_{\mathrm{levels}}
  :=
  \#\bigl\{k\in\Nzero:\text{there is a }t\in[0,T_*]\text{ with }H(t)=L_k\bigr\}.
\]
The roof leaves room for only finitely many such doublings:
\begin{equation}\label{eq:BodyFiniteZenoOctaveCount}
  N_{\mathrm{levels}}
  \leq
  \begin{cases}
    0,&H_{\max}<H_0,\\
    1+\left\lfloor\log_2(H_{\max}/H_0)\right\rfloor,
      &H_{\max}\geq H_0.
  \end{cases}
\end{equation}
This conclusion never asks how quickly one level is crossed.
The summable heat clocks from Section~\ref{sec:CriticalHeightAndDerivativeTower} cannot
assemble an infinite ladder because the height itself runs out of levels.
A width might still shrink while the height stays bounded, so the geometric
picture has one more question to answer.

The next rectangle, \(\cR_{0,2}\), follows the motion below the roof rather than
only its maximum:
\begin{equation}\label{eq:LowRectangleZeroTwo}
  u\in L^2(I_*;H^3),
  \qquad
  u_t\in L^2(I_*;H^1),
  \qquad
  u\in C([0,T_*];H^2).
\end{equation}
The two adjacent rows now contain both factors in the derivative
\[
  H'(t)
  =\operatorname{Re}
  \pair{\Lambda^{1/2}u(t)}{\Lambda^{1/2}u_t(t)}_{L^2},
\]
and Cauchy--Schwarz pays its total variation:
\begin{equation}\label{eq:CriticalHeightTotalVariationBound}
  \int_0^{T_*}|H'(t)|\,dt
  \leq
  \norm{\Lambda^{1/2}u}_{L^2_tL^2_x}
  \norm{\Lambda^{1/2}u_t}_{L^2_tL^2_x}
  <\infty.
\end{equation}
The same rectangle has \(E_{3/2}\in L^1(I_*)\).  Hence the completed critical row
\[
  \mathcal P_{1/2}=H'+2\nu E_{3/2}
\]
puts the signed production in \(L^1(I_*)\), and integration records exactly how
much of it has accumulated:
\begin{equation}\label{eq:BodySignedCriticalPrefix}
  \int_0^t\mathcal P_{1/2}(s)\,ds
  =H(t)-H(0)+2\nu\int_0^tE_{3/2}(s)\,ds.
\end{equation}
The two terms on the right are already finite, so
\begin{equation}\label{eq:BodySignedCriticalPrefixBound}
  A_*(T_*)
  :=\sup_{0<t<T_*}\int_0^t\mathcal P_{1/2}(s)\,ds
  <\infty.
\end{equation}
Thus the signed prefix is read off after the rectangle has been built; it is not
an upstream estimate from which the rectangle was obtained.
Sharp rises and falls remain possible, but their net signed accumulation cannot
diverge before \(T_*\).

For a vortex, the decisive entry in \(\cR_{0,2}\) is
\(u\in L^2(I_*;H^3)\).  In three dimensions
\(H^3\hookrightarrow W^{1,\infty}\), whence
\[
  \int_0^{T_*}\norm{\nabla u(t)}_\infty\,dt
  \leq
  T_*^{1/2}\norm{u}_{L^2(I_*;H^3)}<\infty.
\]
Both vorticity and strain are parts of this gradient.  Their accumulated action
before \(T_*\) is finite in the precise sense
\begin{equation}\label{eq:FiniteLipschitzVorticityStrainAction}
  \int_0^{T_*}
  \left(
    \norm{\nabla u(t)}_\infty
    +\norm{\omega(t)}_\infty
    +\norm{S(t)}_\infty
  \right)dt
  <\infty.
\end{equation}

\subsection{The two vortex readings}\label{sub:ClosingCriticalHeightZenoAndTheVortex}

What does a narrowing core mean for the fluid inside it?
There are two answers to separate.  The first follows the same material
particles.  The second allows the visible concentration to replace its
particles as it moves.  Neither exhausts singular geometry, but each makes a
different claim that the rectangle bounds can test.

Follow the fixed labels first.
Let \(X(a,t)\) solve
\[
  \partial_tX(a,t)=u(X(a,t),t),
  \qquad
  X(a,0)=a.
\]
Two such trajectories satisfy
\[
  \frac d{dt}|X(a,t)-X(b,t)|
  \geq
  -\norm{\nabla u(t)}_\infty|X(a,t)-X(b,t)|,
\]
and the finite Lipschitz action integrates this differential inequality to
\begin{equation}\label{eq:BodyMaterialCoreLowerBound}
  |X(a,t)-X(b,t)|
  \geq
  |a-b|
  \exp\!\left(
    -\int_0^t\norm{\nabla u(s)}_\infty\,ds
  \right)
  >0
\end{equation}
Because the same row also puts \(u\) in \(L^1(I_*;L^\infty)\), every trajectory
has a limit at \(T_*\).  Letting \(t\uparrow T_*\) preserves the lower bound:
\[
  |X(a,T_*)-X(b,T_*)|
  \geq |a-b|\exp\!\left(-\int_0^{T_*}\norm{\nabla u(s)}_\infty\,ds\right)>0.
\]
Distinct labels remain spatially distinct at the terminal time.
If a drawn core collapses to one point, it cannot still contain exactly the
material core with which it began.

Even perfect centring cannot rescue a nonzero circulation from this fixed-label
picture.  The finite coordinate \((r,m)=(0,3)\) gives
\begin{equation}\label{eq:BodyUniformLipschitzBound}
  L_*:=\sup_{0\leq t\leq T_*}\norm{\nabla u(t)}_\infty<\infty.
\end{equation}
Relative velocity at any chosen centre \(x_c(t)\) obeys
\[
  |u(t,x)-u(t,x_c)|\leq L_*|x-x_c|.
\]
Around a circular cross-section \(C_r(t)\), the constant centre velocity
integrates to zero.  What remains is bounded by its Lipschitz variation:
\begin{equation}\label{eq:BodyCenteredCirculationBound}
  \left|
    \oint_{C_r(t)}u(t,x)\cdot d\ell
  \right|
  \leq2\pi L_*r^2
  \longrightarrow0
  \qquad(r\downarrow0).
\end{equation}
As \(r\) shrinks, the tangential velocity relative to the centre is \(O(r)\)
and the circuit has length \(O(r)\).  Their product is \(O(r^2)\), so this
centred material core cannot carry a nonzero limiting circulation.

Now let membership change.  Particle separation no longer tracks the visible
core, but repeated recruitment would have to keep transferring critical content
to ever smaller wavelengths.  That claim belongs to the high-frequency tail of
the whole velocity field.
For \(s>1/2\) and \(K\geq1\), define
\[
  H_{>K}(t)
  :=\frac12\int_{|\xi|>K}
  |\xi|\,|\widehat u(t,\xi)|^2\,d\xi.
\]
For every \(s>1/2\), the tail satisfies
\begin{equation}\label{eq:BodyUltravioletCriticalTail}
  H_{>K}(t)
  \leq
  \frac12K^{1-2s}\norm{u(t)}_{\dot H^s}^2.
\end{equation}
Choose one integer \(m\geq s\).  The compatible assembly and adjacent trace
proved in \Cref{prop:BodyCompatibleProjectiveAssembly,lem:AdjacentHilbertScaleTrace},
followed by \(H^m\hookrightarrow\dot H^s\), make the estimate uniform in time
and force
\begin{equation}\label{eq:BodyUniformUltravioletVanishing}
  \lim_{K\to\infty}
  \sup_{0\leq t\leq T_*}H_{>K}(t)=0.
\end{equation}
Recruitment may move activity between scales, but it cannot leave a fixed amount
of critical height beyond every cutoff.
The two readings fail for different reasons: fixed labels cannot collapse and
their centred circulation vanishes; changing labels cannot sustain the required
ultraviolet critical tail.

\subsection{The singularity conditions fail at both scales}\label{sub:TheLocalSingularityPictureCloses}

The first rectangle already reaches the global signature of a singular time.
The endpoint theorem of Escauriaza, Seregin, and \v{S}ver\'ak says that a finite
maximal time forces the critical \(L^3\) norm to escape
\parencite{EscauriazaSereginSverak2003}, as recorded in
\eqref{eq:FiniteTimeLThreeEscape}.  The same velocity has the opposite bound in
\eqref{eq:EndpointCriterionFromLowRectangle}.  A finite terminal time cannot
satisfy both statements.

The local signature disappears from the same intersection.  Fix a finite
\(T>0\).  The mixed velocity bounds and the normalized pressure recurrence in
\Cref{prop:PressureCompleteEndpointJet} give two finite coordinates with the
common ceiling
\begin{equation}\label{eq:FiniteIntervalVelocityPressureSupremum}
  K_T
  :=\sup_{0\leq t\leq T}
  \left(
    \norm{u(t)}_{L^\infty}
    +\norm{p(t)}_{L^\infty}
  \right)
  <\infty.
\end{equation}
Now place a parabolic cylinder around any point of the history:
\[
  Q_\rho(z_0)
  =B_\rho(x_0)\times(t_0-\rho^2,t_0)
  \subset\R^3\times(0,T].
\]
Its spatial volume and time length contribute \(\rho^5\).  Dividing by the
critical factor \(\rho^2\) still leaves three powers of the radius:
\begin{equation}\label{eq:VanishingLocalScaleInvariantContent}
  \rho^{-2}
  \iint_{Q_\rho(z_0)}
  \left(|u|^3+|p|^{3/2}\right)\,dx\,dt
  \leq
  C\rho^3\left(K_T^3+K_T^{3/2}\right)
  \longrightarrow0
  \qquad(\rho\downarrow0).
\end{equation}
A Caffarelli--Kohn--Nirenberg singular point would require a positive
scale-invariant amount of velocity-pressure content to survive along shrinking
cylinders \parencite{CaffarelliKohnNirenberg1982}.  Here that content vanishes.

This is the return promised by the vortex picture.  The alleged endpoint has
become a smooth time slice.  The critical height has a finite roof, so the Zeno
ladder has only finitely many levels.  The stretching action is integrable, a
material core cannot collapse its labels, and the local velocity-pressure
content loses the mass a singular point would need at every smaller scale.

\clearpage
\section{Further Consequences and the Euler Boundary}\label{sec:FurtherConsequencesAndEulerBoundary}

The four-part argument is complete once the intersection has returned to the
critical height, the Zeno ladder, and the singularity conditions.  Two further
projections remain useful.  One adds the vorticity across all dyadic scales;
the other removes viscosity and shows which rung of the Tower then disappears.

\subsection{A finite sum across infinitely many scales}\label{sub:AFiniteSumAcrossScales}

Fix a finite \(T>0\).  On \([0,T]\), the next spatial row
\(\cR_{0,2}[u_0;T]\) puts \(u\) in \(L^2_tH^3_x\).  That one extra row is
enough to add the vorticity from every dyadic wavelength rather than inspect
the shells one at a time.

Let \((\Delta_j)_{j\geq-1}\) be an inhomogeneous Littlewood--Paley
decomposition \parencite[Chapter~2]{BahouriCheminDanchin2011}.  At high
frequency, \(\Delta_j\omega\) sees wavelengths of order \(2^{-j}\), while its
\(L^\infty\) norm records the largest rotation in that shell.  Define
\[
  \norm{\omega}_{B^0_{\infty,1}}
  :=\sum_{j\geq-1}\norm{\Delta_j\omega}_\infty.
\]
Bernstein's inequality costs three halves of a derivative in passing from
\(L^2\) to \(L^\infty\):
\[
  \norm{\Delta_j\omega}_\infty
  \leq C2^{3j/2}\norm{\Delta_j\omega}_2.
\]
The \(H^2\) weight of \(\omega\) supplies two derivatives.  Splitting
\(2^{3j/2}=2^{-j/2}2^{2j}\) exposes the spare half derivative, and
Cauchy--Schwarz sums it:
\begin{align*}
  \sum_{j\geq0}\norm{\Delta_j\omega}_\infty
  &\leq
  C\left(\sum_{j\geq0}2^{-j}\right)^{1/2}
  \left(\sum_{j\geq0}2^{4j}\norm{\Delta_j\omega}_2^2\right)^{1/2}
  \\
  &\leq C\norm{\omega}_{H^2}
  \leq C\norm{u}_{H^3}.
\end{align*}
The low shell is controlled directly by \(\norm{\omega}_2\).  Thus one finite
Sobolev coordinate pays the full \(\ell^1\) shell sum:
\begin{equation}\label{eq:VorticityBesovFromHThree}
  \norm{\omega}_{B^0_{\infty,1}}
  \leq C\norm{u}_{H^3}.
\end{equation}
Cauchy--Schwarz in time now turns the \(L^2_tH^3_x\) row into
\begin{equation}\label{eq:FiniteDyadicVorticityAction}
  \sum_{j\geq-1}
  \int_0^{T}\norm{\Delta_j\omega(t)}_\infty\,dt
  =\int_0^{T}\norm{\omega(t)}_{B^0_{\infty,1}}\,dt
  <\infty.
\end{equation}
Infinitely many shells may still become active.  What cannot happen before this
finite \(T\) is a nonsummable succession of their time-integrated peaks.
The conclusion consumes only \(\cR_{0,2}\): it asserts neither an analytic
radius nor a weighted estimate over all derivative orders.

\Needspace{8\baselineskip}
\subsection{Why the Euler Tower is different}\label{sub:WhyTheEulerTowerIsDifferent}

Remove viscosity while keeping the two histories separate.  Both equations can
still be differentiated in time, so both possess a mixed jet.  Write
\[
  V_n^{\mathrm{NS}}:=\partial_t^nu^{\mathrm{NS}},
  \qquad
  Q_n^{\mathrm{NS}}:=\partial_t^np^{\mathrm{NS}},
\]
and define \(V_n^{\mathrm E},Q_n^{\mathrm E}\) from an Euler solution with its
own history.  The two recurrences differ in one place:
\begin{align}
  V_{n+1}^{\mathrm{NS}}
  &+
  \sum_{a=0}^{n}\binom{n}{a}
  (V_a^{\mathrm{NS}}\cdot\nabla)V_{n-a}^{\mathrm{NS}}
  +\nabla Q_n^{\mathrm{NS}}
  =\nu\Delta V_n^{\mathrm{NS}},
  \label{eq:NSTemporalRecurrenceComparison}
  \\
  V_{n+1}^{\mathrm E}
  &+
  \sum_{a=0}^{n}\binom{n}{a}
  (V_a^{\mathrm E}\cdot\nabla)V_{n-a}^{\mathrm E}
  +\nabla Q_n^{\mathrm E}
  =0.
  \label{eq:EulerTemporalRecurrenceComparison}
\end{align}
Transport and the nonlocal pressure response remain in both equations.
Only the Navier--Stokes row contains \(\nu\Delta V_n^{\mathrm{NS}}\), which
connects a temporal response to two more spatial derivatives of that same
response.

The energy rows expose the same difference.  Set
\[
  E_s^{\mathrm{NS}}:=\frac12\norm{\Lambda^su^{\mathrm{NS}}}_2^2,
  \qquad
  E_s^{\mathrm E}:=\frac12\norm{\Lambda^su^{\mathrm E}}_2^2,
\]
with the corresponding nonlinear transfer terms.  Then
\begin{equation}\label{eq:NSEulerEnergyComparison}
  (E_s^{\mathrm{NS}})'
  =\mathcal P_s^{\mathrm{NS}}-2\nu E_{s+1}^{\mathrm{NS}},
  \qquad
  (E_s^{\mathrm E})'
  =\mathcal P_s^{\mathrm E}.
\end{equation}
Repeated time differentiation of the Navier--Stokes row, followed by division
by \((-2\nu)^n\), isolates \(E_{n+1/2}^{\mathrm{NS}}\) in
\eqref{eq:BodyCompletedCriticalRow}.  The Euler row still mixes time and space
nonlinearly, but no signed coercive term singles out the next spatial height.

If an Euler history were independently known to lie in the analogous
all-orders intersection, then every fixed Sobolev norm would be controlled and
in particular it would satisfy the Beale--Kato--Majda condition
\(\int_0^T\|\omega^{\mathrm E}(t)\|_\infty\,dt<\infty\)
\parencite{BealeKatoMajda1984}.  The present theorem supplies no such Euler
intersection, and none of its constants is asserted to remain bounded as
\(\nu\downarrow0\).

That is where the argument ends.  Euler keeps the jet, transport, pressure, and
their nonlinear coupling; it loses the viscous rung that converts temporal
response into a separately controllable spatial ascent.  At every fixed
\(\nu>0\), that ascent exposes the finite rectangles whose terminal bounds rule
out a last time.  The resulting Navier--Stokes history has every finite mixed
coordinate on each finite interval, while its kinetic energy remains uniformly
bounded for all time.
